\section{Research Gap}
\label{sec_research_gap}

The comparative analysis in Section~\ref{sec_comparative_analysis} establishes that autonomous VAPT systems have so far addressed at most one of three capabilities at a time. Table~\ref{tab:gap-table} consolidates this finding across the reviewed literature.

\begin{table}[htbp]
	\centering
	\small
	\caption{Three-dimensional gap in the reviewed autonomous VAPT literature: what each capability area shows and where it falls short.}
	\label{tab:gap-table}
	\begin{tabular}{@{}p{3.2cm}p{5.5cm}p{5.2cm}@{}}
		\toprule
		\textbf{Capability area}            & \textbf{What the reviewed literature shows}                                                                                                                                     & \textbf{Documented limitation}                                                                                                                           \\
		\midrule
		Exploration breadth                 & HPTSA and VulnBot scale to wide, unfamiliar attack surfaces through hierarchical team dispatch and phase-sequenced PTG planning                                                 & Neither models prerequisite relationships between discovered vulnerabilities; no formal notion of which weakness must precede another                    \\
		Prerequisite / dependency reasoning & CHECKMATE introduces PDDL-style explicit prerequisite modeling; PentestEval's GT-ADM ablation shows that expert-quality ADM raises end-to-end success from 0.31 to 0.67 (+0.14) & Both require weaknesses to be pre-enumerated by the operator before planning begins; neither constructs a dependency model from dynamic discovery output \\
		Cross-session memory                & No system in the reviewed set retains or reuses strategies across engagements against the same or similar targets                                                               & An open area this thesis treats as secondary to the above two gaps; empirical value to be assessed once the primary gaps are addressed                   \\
		Cross-surface evaluability          & Every autonomous system in the reviewed set evaluates on a single attack-surface family                                                                                         & Architectural decisions cannot be distinguished from surface-specific tuning effects in any existing study                                               \\
		\bottomrule
	\end{tabular}
\end{table}

HPTSA and VulnBot both search wide attack surfaces, but neither carries an explicit model of which candidate weaknesses must precede others. When a phase produces multiple candidates, they are dispatched without an ordering that reflects prerequisite structure. CHECKMATE fills precisely this gap --- its PDDL-style planning operators encode prerequisite relationships explicitly --- but at a cost: relevant weaknesses must be enumerated by the operator before the planning phase begins, and the system is not designed for open-ended discovery on an unfamiliar surface. PrediQL occupies a third position: its dependency graph is schema-derived, provided by the GraphQL API itself rather than assembled from the agent's own discovery output, and the approach is surface-specific to GraphQL by design.

The two benchmark studies in this selection quantify the cost of each missing capability. CVE-Bench's failure-mode analysis finds that insufficient exploration accounts for 37.5\%--80.0\% of unsuccessful attempts across every agent and evaluation setting; its best-performing agent achieves only 13\% pass@5 \citep{zhu2025cvebench}. PentestEval's stage-level breakdown finds that attack decision-making --- the stage that requires reasoning about which discovered weakness to pursue in which order --- achieves the lowest average performance of any planning stage (Spearman correlation 0.25 across nine evaluated LLMs) and delivers the largest single-stage improvement when replaced with expert ground truth (+0.14, raising end-to-end success from 0.31 to 0.67) \citep{yang2025pentesteval}. The two failure modes are not independent: a system that searches widely without a prerequisite model will waste iterations on unreachable paths; a system that reasons about dependencies only over a pre-supplied weakness set cannot operate on a surface it has never seen. Addressing one without the other leaves the compound gap intact.

The gap this thesis investigates is therefore precise: no reviewed system combines (i)~open-ended exploration of an unfamiliar attack surface with (ii)~a dynamically constructed, prerequisite-aware model of dependencies between discovered weaknesses, and (iii)~evaluation across more than one independently benchmarked attack-surface family. Chapter~\ref{chapter_proposedMethod} introduces the proposed RedGrid architecture as a first attempt to address this compound gap.
